\section{Implicaciones organizacionales y de inversión}

\subsection{Estructura de talento y cultura}

Burak enfatiza en particular la importancia de la densidad de talento y de una cultura de iteración: organizar los squad multifuncionales en una tríada de product + safety + research. Dentro de cada squad se necesitan cuatro tipos de roles —prompt engineer, search engineer, ops engineer y policy analyst— para asegurar que el agent pipeline no solo pueda ejecutarse, sino también administrarse.

La experiencia de Anthropic en 2023 es que la densidad de talento importa más que la cantidad. Definieron el valor organizacional como «quién puede ejecutar de inmediato un experimento y explicar el failure». Esta cultura asegura que el code review cubra a la vez quality y compliance.

\subsection{Asignación de capital y portafolio de proyectos}

La inversión en GenAI no se apuesta enteramente en un solo agent, sino que se construye un portfolio de varias etapas:

\begin{itemize}
    \item Stage 1: GPT-style capability research continúa explorando scaling, architecture
    \item Stage 2: Platform / Agent ops brinda tooling + governance
    \item Stage 3: Conversión en producto convierte el agent en specific vertical solutions (como finance assistant)
\end{itemize}

\begin{importantbox}{Beneficios del portfolio por etapas}
Esta combinación permite que el equipo explore la innovación en el early-stage, garantice platform reliability en el mid-stage y asuma la commercialization en el late-stage. Los metrics / budget de cada etapa deben mantenerse claramente separados, para evitar que el departamento de capacidades y el departamento de product interfieran entre sí dentro del mismo sprint.
\end{importantbox}

\subsection{Resumen de la sección}

En el nivel organizacional es necesario establecer tres mecanismos clave —«densidad de talento», «squad multifuncionales» y «inversión por etapas»— para mantener al equipo bajo control mientras la tecnología evoluciona con rapidez.

